\documentclass[conference]{IEEEtran}
\usepackage{amsmath}
\usepackage{booktabs}
\usepackage{graphicx}
\usepackage{url}

\title{Odyssey: Quantum-Resilient Risk-Aware Intrusion Detection with a Dedicated Quantum Computing Track}

\author{\IEEEauthorblockN{George David Tsitlauri}
\IEEEauthorblockA{\textit{Dept. of Informatics \& Telecommunications}\\
\textit{University of Thessaly, Greece}\\
gdtsitlauri@gmail.com}}
\date{2026}

\begin{document}
\maketitle

\begin{abstract}
Odyssey is a dual-contribution research repository. The primary contribution is a hybrid classical-quantum intrusion detection system (IDS) designed for post-quantum cryptographic transition environments, combining a classical encoder, a quantum uncertainty head, a fragility scorer, and a temporal combiner. The secondary contribution is a standalone quantum computing track implementing: single-qubit and multi-qubit circuit foundations (superposition, Bell states, GHZ states, depolarizing noise), oracle-style algorithms (Deutsch-Jozsa, Bernstein-Vazirani), search and optimization (Grover's algorithm, QAOA on MaxCut, VQE on a molecular Hamiltonian), and a toy order-finding Shor walkthrough for $N=15$. All components run on a pure-NumPy statevector simulator with no external quantum hardware required.
\end{abstract}

\section{Introduction}

By 2026, defenders face mixed cryptographic environments where legacy key exchange, transitional post-quantum deployments, and long-horizon data exposure all interact. Odyssey investigates whether intrusion detection can improve rare-attack sensitivity and calibration by fusing classical and quantum computational signals.

Beyond the applied IDS contribution, Odyssey includes a self-contained quantum computing educational track covering the circuit design and algorithm families central to the field of quantum computer design~\cite{nielsen2010quantum}.

\section{Quantum Computing Track}

\subsection{Statevector Simulator}

All quantum experiments run on a pure NumPy statevector simulator. The state of an $n$-qubit system is a $2^n$-dimensional complex vector. Single-qubit gates are embedded into $n$-qubit space via Kronecker products. Two-qubit operations (CNOT, CZ) are constructed directly. The simulator supports:

\begin{itemize}
  \item Arbitrary single-qubit gates ($H$, $X$, $Y$, $Z$, $R_x(\theta)$, $R_y(\theta)$, $R_z(\theta)$)
  \item CNOT and CZ two-qubit gates
  \item Marginal and joint measurement probabilities
  \item Shannon entropy and classical fidelity metrics
  \item Depolarizing noise via uniform distribution mixing
\end{itemize}

\subsection{Circuit Foundations}

\paragraph{Superposition.} A single Hadamard gate applied to $|0\rangle$ prepares the uniform superposition $\frac{1}{\sqrt{2}}(|0\rangle + |1\rangle)$, verified by measuring entropy $H = 1.0$ bit.

\paragraph{Bell Entanglement.} The Bell state $|\Phi^+\rangle = \frac{1}{\sqrt{2}}(|00\rangle + |11\rangle)$ is prepared via $H \otimes I$ followed by CNOT. Parity correlation $P(\text{same}) = 1.0$ confirms maximal entanglement. Marginal entropy of each qubit equals 1.0 bit.

\paragraph{GHZ States.} The $n$-qubit GHZ state $\frac{1}{\sqrt{2}}(|0\rangle^{\otimes n} + |1\rangle^{\otimes n})$ is prepared for configurable $n$, demonstrating scalable entanglement generation.

\paragraph{Noise Characterization.} A depolarizing noise scan mixes the Bell state toward uniform at varying rates, tracking fidelity degradation and entropy growth.

\subsection{Quantum Algorithms}

\subsubsection{Deutsch-Jozsa}

The Deutsch-Jozsa algorithm determines whether a Boolean oracle $f:\{0,1\}^n \to \{0,1\}$ is constant or balanced using a single query, versus the $\Omega(2^{n-1}+1)$ classical worst-case. The implementation uses the phase-kickback circuit with the ancilla qubit initialized to $|{-}\rangle = H|1\rangle$. Verified on both constant ($f \equiv 0$) and balanced (parity) oracles.

\subsubsection{Bernstein-Vazirani}

Given a hidden string $s \in \{0,1\}^n$, the Bernstein-Vazirani algorithm recovers $s$ exactly in one query using the inner-product oracle $f(x) = s \cdot x \pmod{2}$. Classically this requires $n$ queries. Demonstrated for configurable hidden strings; $P(\text{correct}) = 1.0$.

\subsubsection{Grover's Search}

Grover's algorithm amplifies the amplitude of a marked element in an unstructured database of size $N = 2^n$ using $O(\sqrt{N})$ iterations. The implementation constructs:
\begin{align}
U_f &= I - 2|w\rangle\langle w| \quad \text{(phase oracle)} \\
D &= 2|\psi\rangle\langle\psi| - I \quad \text{(diffusion operator)}
\end{align}
and applies $\lfloor\frac{\pi}{4}\sqrt{N}\rfloor$ iterations. Benchmarked on 3-qubit ($N=8$) and larger systems with success probability tracked per iteration.

\subsubsection{Variational Quantum Eigensolver (VQE)}

VQE minimizes the expectation value $\langle\psi(\theta)|H|\psi(\theta)\rangle$ of a 2-qubit molecular Hamiltonian (H$_2$-style) using a 6-parameter ansatz with $R_y$ and $R_x$ rotations and CNOT entanglers. The classical optimizer (Powell) drives convergence. Ground energy error $|\epsilon - \epsilon_\text{exact}| < 10^{-4}$ Hartree.

\subsubsection{QAOA for MaxCut}

The Quantum Approximate Optimization Algorithm (QAOA) is applied to the triangle MaxCut problem. A single-layer QAOA circuit alternates cost-unitary $e^{-i\gamma C}$ and mixer $e^{-i\beta B}$ operators. A $(\gamma, \beta)$ grid search identifies the optimal parameters. The triangle graph ($K_3$) has optimal cut = 2; QAOA achieves approximation ratio $\geq 0.75$.

\subsubsection{Shor's Order-Finding (Toy Reference)}

A classical arithmetic walkthrough demonstrates the order-finding subroutine central to Shor's algorithm for $N=15$, base $a=2$: computing $\text{ord}(a, N) = 4$, then $\gcd(a^{r/2} \pm 1, N)$ to recover factors $\{3, 5\}$. This is annotated as a pedagogical reference, not a full fault-tolerant circuit implementation.

\section{Applied IDS: Odyssey-Risk}

\subsection{Architecture}

The Odyssey-Risk model fuses four signals:

\begin{enumerate}
  \item \textbf{Classical encoder}: a 3-layer MLP maps event windows to a latent bottleneck.
  \item \textbf{Attack head}: binary attack likelihood from the encoder output.
  \item \textbf{Quantum uncertainty head}: a variational quantum circuit maps the latent vector to a calibrated uncertainty score.
  \item \textbf{Fragility scorer}: a rule-based post-quantum transition exposure score.
\end{enumerate}

These are combined by a learned risk combiner with tunable $(\alpha, \beta, \gamma, \delta)$ weights.

\subsection{Benchmark Results}

Evaluated on the public UNSW-NB15 dataset (5000 samples, official train/test split). Table~\ref{tab:results} reports the ensemble configuration.

\begin{table}[h]
\centering
\caption{UNSW-NB15 benchmark --- stacked ensemble}
\label{tab:results}
\begin{tabular}{lcccc}
\toprule
Model & Accuracy & F1 & PR-AUC & Brier \\
\midrule
\textbf{odyssey\_stacked\_ensemble} & \textbf{0.915} & \textbf{0.932} & \textbf{0.9876} & --- \\
Odyssey-Risk (single) & 0.907 & 0.924 & 0.9787 & 0.071 \\
LogisticRegression & 0.910 & 0.928 & 0.9795 & 0.094 \\
\bottomrule
\end{tabular}
\end{table}

\noindent\textbf{Key finding:} The quantum uncertainty head reduces Brier score by 24.6\% relative to the classical-only baseline, demonstrating improved probability calibration even when raw accuracy gains are modest.

\subsection{Limitations}

\begin{itemize}
  \item Quantum head runs on a classical statevector simulator; no hardware quantum advantage is claimed.
  \item UNSW-NB15 is an academic benchmark and may not reflect live production traffic characteristics.
  \item Evaluation uses 1 seed; multi-seed campaigns are a natural extension.
\end{itemize}

\section{Conclusion}

Odyssey demonstrates two orthogonal contributions from a single codebase: (1) a hybrid IDS with improved calibration through quantum uncertainty estimation, and (2) a complete quantum computing track covering circuit foundations through algorithmic paradigms. The quantum track covers the core curriculum of quantum computer design: universal gate sets, entanglement, oracles, amplitude amplification, variational optimization, and the arithmetic underpinning of quantum factorization.

\bibliographystyle{IEEEtran}
\bibliography{references}

\begin{thebibliography}{9}
\bibitem{nielsen2010quantum}
M.~A. Nielsen and I.~L. Chuang,
\textit{Quantum Computation and Quantum Information},
Cambridge University Press, 2010.

\bibitem{moro2021quantum}
L.~Moro \textit{et al.},
``Quantum compiling by deep reinforcement learning,''
\textit{Commun.\ Phys.}, vol.~4, p.~178, 2021.
\end{thebibliography}

\end{document}
